\scanpage{3}
\mspage{- 111 -}
% archive sheet number: 2
% lines 1–6
irréductible de $f^{-1}(y)$ \deleted{\illegible{2}}, de point générique $x$, est contenue dans l'adhérence d'un voisinage ouvert affine de $x$. On peut aussi remplacer $Y$ par le schéma local de $y$, donc supposer que \deleted{\illegible{1}} $Y=\mathrm{Spec}(B)$, où $B$ est un anneau local \inserted{noethérien}, d'idéal maximal $\handwritten{\underline}{m}==\handwritten{\underline}{j}_y$, et $X=\mathrm{Spec}(A)$, où $A$ est une algèbre \deleted{\illegible{2}} \inserted{noethérienne} sur $B$ ; on peut en outre (remplaçant $B$ par son image dans $A$) supposer $B\handwritten{\subset} A$.
% slanted
\margin{\handwritten{préparatifs laborieux inutiles}}

% lines 6–11
Alors il s'agit de voir que si $\handwritten{\underline}{p}$ est \deleted{\illegible{1}} un idéal premier\deleted{\unsure{m}} minimal associé à $A\handwritten{\underline}{m}$, on a $\mathrm{rg}(\handwritten{\underline}{p})\handwritten{\leq}\mathrm{rg}(\handwritten{\underline}{m})$. Or, dans l'anneau local \deleted{$\unsure{A|{}}_{\unsure{p}}$} $A_{\handwritten{\underline}{p}}$ l'idéal maximal $\handwritten{\underline}{p}A_{\handwritten{\underline}{p}}$ est engendré par $A\handwritten{\underline}{m}$, donc admet un système de $r$ générateurs, si $r$ est le rang de $\handwritten{\underline}{m}$, ce qui démontre notre assertion.

% lines 12–15
\uline{Corollaire 1 .- Soient} \inserted{$X$ \uline{et}} $Y$ \uline{\deleted{\unsure{un}} \inserted{deux} préschémas localement noethériens}, $f:X\handwritten{\to Y}$ \uline{un morphisme \deleted{\illegible{2}}}, $Y'$ \uline{une partie fermée irréductible de} $Y$, $X'$ \uline{une composante irréductible de} $f^{-1}(Y')$, \uline{dont l'image dans} $Y$ \uline{soi\handwritten{t} dense dans} $Y'$. \uline{Alors} $\mathrm{codim}_X(X')\handwritten{\leq}\mathrm{codim}_Y(Y')$.

% lines 16–18
Si $x$ et $y$ sont les points génériques de $X'$ et $Y'$ respectivement, \deleted{\illegible{1}} l'hypothèse sign\unsure{i|j}fie que $f(x)=y$, et le cor.1 n'est qu'une autre fo\handwritten{r}me de la prop.8 .

% lines 19–22
\uline{Corollaire 2 .- Soient} $X$ \uline{un préschéma \deleted{\unsure{noethérien}} \inserted{localement noethérien},} $f$ \uline{une section de} $O_X$ \uline{au-dessus de} $X$, \uline{qui n'est identiquement nulle dans aucune composante irréductible de} $X$. \uline{Soit} $Z_f$ \uline{l'ensemble des} $x\handwritten{\in}X$ \uline{tels que} $f(x)\handwritten{=0}$ \uline{Alors toute composante irréductible de} $Z_f$ \uline{est de codimension} 1 .

% lines 23–29
Il suffit de démontrer qu'il en est ainsi dans un vois\unsure{i}nage ouvert affine de tout point de $Z_f$, donc on peut supposer \deleted{\illegible{1}} $X=\mathrm{Spec}(A)$ affine d'anneau $A$ noethérien ; par hypothèse $f$ n'est \deleted{\illegible{1}} contenu dans aucun des idéaux premiers minimaux de $A$, et tous les idéaux premier\handwritten{s} minimaux associés à l'idéal principal $(f)$ sont donc de rang 1 (``Hauptidealsatz'' de Krull ;cf.Bourbaki,Alg.comm.\qquad), ce qui prouve le corollaire. Ce dernier se déduit aussi de la prop.8, en \deleted{\illegible{1}} consi-
% slanted
\margin{\handwritten{affinisation inutile, si on remarque que pour que $x\in X$ soit \deleted{le} \unsure{pt} générique d'une comp. irréd. de $Z_f$, il faut et \unsure{su.} que $f\unsure{\in|+}\underline O_x$ \deleted{\unsure{engendre}} un \deleted{\unsure{x|l'}}idéal de définition}}
